% =====================================================================
% Section 6 — Subsystem Benchmark Roll-Up
% Fragment \input inside \section{Subsystem Benchmark Roll-Up}.
% Numbers per _DRAFTING_BRIEF.md section 4 + VALIDATION_MATRIX.md claim map.
% =====================================================================

\subsection{Verification and Validation Methodology}
\label{subsec:vv-method}

The model suite was verified and validated under the discipline of the AIAA
Editorial Policy on Numerical and Experimental Accuracy~\cite{aiaa_numerical_policy},
which applies to engineering prediction methods as it applies to computational
fluid dynamics: code verification against analytical or tabulated solutions,
error quantification with an explicit metric and stated regime, and an
iterative-convergence demonstration wherever an iterative solver appears. We
distinguish two classes of evidence. \emph{Code verification} anchors a
subsystem against an exact or tabulated reference---NACA Report~1135 shock and
expansion tables~\cite{naca1135}, the Taylor--Maccoll cone solution~\cite{taylormaccoll1933,anderson2006},
and the Rayleigh pitot relation---where the acceptance criterion is agreement to
near machine precision. \emph{Physical validation} anchors a subsystem against
published wind-tunnel, ballistic-range, free-flight, or CFD data---NACA
RM~A52H28~\cite{a52h28}, NACA TN~3393~\cite{chapman1955}, NASA TM~X-653~\cite{nielsen1962},
and the DTIC AD0487365 cone-cylinder range tests~\cite{grabow1965}---where the
acceptance criterion is a stated mean absolute percent error (MAPE) or mean
absolute error (MAE) gate appropriate to a semi-empirical engineering method.
Iterative solvers (the oblique-shock and Taylor--Maccoll root finders) converge
to a $10^{-12}$ relative tolerance; integrated time-step convergence is reported
in Sec.~\ref{sec:corpus} ($\lvert s\rvert \approx 0.98\%$ per $10\%$ contraction
of the step).

To grade the evidence consistently we adopt a four-level rubric. A row is
\emph{A-level} when it is matched against published external or tabulated data
with a quantitative acceptance criterion; \emph{B-level} when it is a
source-anchored analytical or flight-corpus closure that has not been isolated
against a dedicated published component dataset; \emph{C-level} for internal
consistency or numerical-integrity checks; and \emph{D-level} for a calibrated
heuristic without external closure. Only A-level rows are claimed as validated
in the sense of the AIAA policy; B-level rows are disclosed as such throughout.

\subsection{Externally Benchmarked A-Level Subsystems}
\label{subsec:a-level}

Twenty subsystems currently meet the A-level standard: an independent
external benchmark exists, and the model passes a quantitative acceptance gate
against it. They are collected in Table~\ref{tab:benchmarks},
with the primary reference, Mach range, and headline metric for each. The
tabulated rows span $M = 0.6$ to $M = 10$ across the analytical, atmospheric,
drag, and stability families; corroborating sources for each subsystem are
documented in the per-benchmark validation memos under
\texttt{paper/data/md/}. A single externally anchored \emph{negative} benchmark
(the NACA RM-10 high-fineness finned body) is retained outside the A-level count
to bound and explicitly exclude a geometry family that lies outside the
calibration envelope; it is discussed in Sec.~\ref{sec:limitations} rather than
presented as a validating result.

The analytical and atmospheric foundations verify to near machine precision. The
exported speed-of-sound table agrees with the U.S.\ Standard Atmosphere
1976~\cite{nasa1976ussa} to a maximum error of $0.016\%$; the oblique-shock
$\theta$--$\beta$--$M$ solver reproduces NACA Report~1135 to a maximum angle
error of $0.021\%$; the Prandtl--Meyer expansion matches to a maximum absolute
turning-angle error of $0.004^{\circ}$; and the Taylor--Maccoll cone solver
returns the published cone-shock angle to a maximum relative error of $0.825\%$
against a $1\%$ gate. These are code-verification results: the agreement reflects
correct implementation of an exact theory, not an empirical calibration.

The drag and stability families are physical validations and carry the larger,
honestly reported residuals characteristic of semi-empirical engineering
closures. Nose-and-body foredrag matches NACA RM~A52H28 across five shapes to an
aggregate MAE of $0.029$ (gate $0.035$). Turbulent base drag matches NACA
TN~3393 to a MAPE of $15.9\%$; the laminar base-drag closure attributed to
Chapman matches the corresponding TN~3393 laminar rows to $4.4\%$. Finned-vehicle
total drag, exercised on the Army--Navy Basic Finner reference projectile against
the ballistic-range data of Dupuis and Hathaway~\cite{dupuis1997}, achieves a
MAPE of $11.8\%$ over eight points spanning $M = 1.08$ to $4.30$
(Fig.~\ref{fig:basic-finner}); four pointwise residuals exceed $14\%$, so the
claim is an aggregate-gate pass rather than a uniform pointwise match.
Static stability, validated against the wing-body-tail data of NASA
TM~X-653~\cite{nielsen1962} over $M = 0.6$ to $5.82$, returns a normal-force-slope
MAPE of $6.84\%$ and a center-of-pressure ($x_{CP}/d$) MAPE of $7.11\%$ over ten
points each (Fig.~\ref{fig:x653}). A hypersonic extension of the cone foredrag
method, taken against the DTIC AD0487365 cone-cylinder range
tests~\cite{grabow1965}, is reported as an \emph{exploratory} (B-level) result
rather than a validated A-level subsystem; it is presented with its thin-cone
limitation in Sec.~\ref{subsec:b-level} and is not counted among the A-level
rows above.

\begin{table*}[htbp]
\centering
\caption{Representative externally benchmarked A-level subsystems, with primary
reference, Mach range, and headline metric. Analytical and atmospheric rows are
code-verification results (near-exact agreement); drag and stability rows are
physical validations against published wind-tunnel, ballistic-range, free-flight,
or tabulated data, with the semi-empirical residual reported as a mean absolute
percent error (MAPE) or mean absolute error (MAE).}
\label{tab:benchmarks}
\begingroup\footnotesize
\begin{tabular}{@{}r p{0.215\textwidth} p{0.225\textwidth} p{0.175\textwidth} p{0.235\textwidth}@{}}
\toprule
\# & Subsystem & Primary reference & Mach range & Headline metric \\
\midrule
1  & Speed of sound                & U.S.\ Std.\ Atm.\ 1976~\cite{nasa1976ussa}        & sea level--80\,km          & max error $0.016\%$ \\
2  & Sutherland viscosity          & Incropera / NIST~\cite{incropera2007}             & 150--500\,K                & MAPE $0.54\%$ vs.\ tables \\
3  & Normal-shock relations        & NACA Report~1135~\cite{naca1135}                  & $M = 1.5$--10              & max error $<0.01\%$ \\
4  & Oblique-shock solver          & NACA Report~1135~\cite{naca1135}                  & $M = 1.5$--10, $\theta = 5$--40$^\circ$ & max angle error $0.021\%$ \\
5  & Prandtl--Meyer expansion      & NACA Report~1135~\cite{naca1135}                  & $M = 1.2$--10              & max abs angle error $0.004^{\circ}$ \\
6  & Taylor--Maccoll cone flow     & Cone tables~\cite{taylormaccoll1933,anderson2006} & $M = 1$--10                & shock angle $0.825\%$ (gate $1\%$) \\
7  & Rayleigh pitot $C_{p,\max}$   & NACA Report~1135~\cite{naca1135}                  & $M = 1$--10, 15 pts        & max error $<0.01\%$ \\
8  & ShockGeometry pre-pass state  & NACA 1135 + Taylor--Maccoll                       & $M = 1.5$--5               & cone surface Mach $0.00\%$ \\
9  & Nose/body foredrag, 5 shapes  & NACA RM~A52H28~\cite{a52h28}                      & $M = 1.5$--3, 5 shapes     & aggregate MAE $0.029$ (gate $0.035$) \\
10 & Base drag, turbulent          & NACA TN~3393~\cite{chapman1955}                   & $M = 2.73$--4.48           & MAPE $15.9\%$ \\
11 & Base drag, laminar (Chapman)  & NACA TN~3393 laminar~\cite{chapman1955}           & $M = 2.73$--4.48           & MAPE $4.4\%$ \\
12 & Fin wave drag (DATCOM 4.1.5.1)& NACA TN~3650~\cite{ulmann1956}                    & $M = 1.1$--2.5, 12 pts     & MAPE $21\%$ ($0.00\%$ vs.\ Ackeret) \\
13 & Compressible skin friction (Van Driest~II) & Hopkins \& Inouye~\cite{hopkins1971} & $M = 1.5$--9        & best-of-class transformation \\
14 & Static stability $C_{N_\alpha}$ & NASA TM~X-653~\cite{nielsen1962}                & $M = 0.6$--5.82, 10 pts    & MAPE $6.84\%$ \\
15 & Static stability $x_{CP}/d$   & NASA TM~X-653~\cite{nielsen1962}                  & $M = 0.6$--5.82, 10 pts    & MAPE $7.11\%$ \\
16 & Body-alone pitch damping $C_{m_q}$ (no fins) & Tobak NACA TN~3788~\cite{tobak1956}               & $M = 1.5$--3               & $39\%$ at $M=1.5$; sign-correct, conservative (distinct from the finned Basic Finner $C_{m_q}$, Sec.~\ref{subsec:b-level}) \\
17 & Crossflow body $C_{d,c}$      & Jorgensen TR~R-474~\cite{jorgensen1977}           & supersonic crossflow       & exact match (1.20) \\
18 & Crossflow fin $C_{d,c}$       & Hoerner / Jorgensen~\cite{hoerner1965}            & crossflow                  & 1.42 (within range) \\
19 & Magnus body fraction          & Platou 1965~\cite{platou1965}                     & supersonic                 & 0.30 within $0.30$--$0.80$ \\
20 & Finned-vehicle total drag     & Basic Finner~\cite{dupuis1997}                    & $M = 1.08$--4.30, 8 pts    & MAPE $11.8\%$ \\
\bottomrule
\end{tabular}
\endgroup
\end{table*}

\subsection{Secondary and Internally Calibrated Rows}
\label{subsec:secondary}

Two rows that earlier internal documentation listed among the precision A-level
benchmarks are presented here at their honest evidentiary level rather than as
quantitative validations.

The AGARD Model~B subsonic-to-transonic drag comparison against AEDC-TR-70-100~\cite{anderson1970agardb}
is retained only as a \emph{qualitative} secondary trend benchmark. The model
reproduces the subsonic-to-transonic drag-rise trend over $M = 0.2$ to $1.0$, but
its total-drag MAPE on this configuration is $22.6\%$---a loose qualitative
closure, not a precision match---and it is excluded from the A-level count and
from any quantitative claim. The AGARD-B configuration is used in
Sec.~\ref{sec:cfd} as a published-CFD reference geometry, where the same
qualitative-trend framing applies.

The vortex-induced sideforce coefficient is governed by an internal ramp constant
$K_v = 0.20$ that activates above roughly $\alpha = 20^{\circ}$. No verifiable
external benchmark could be located for this coefficient, so it is presented as
an \emph{internally calibrated} parameter constrained by an order-of-magnitude
range check rather than anchored to a published dataset; the earlier
literature-anchored A-level claim is withdrawn. It is reported for transparency
and so that downstream users may re-calibrate it as suitable high-angle-of-attack
data become available.

\subsection{B-Level Disclosures and In-Sample Calibration}
\label{subsec:b-level}

Several model decisions do not meet the A-level standard and are disclosed here.

The pitch-damping derivative $C_{m_q}$ is augmented in production by a fixed
$3\times$ multiplier and a transonic Gaussian peak ($k = 3.5$ at $M \approx 1$),
selected to close the angular-rate residual on the flight corpus. Three
independent reference datasets quantify the consequence, and because they are
distinct data they are reported separately rather than reconciled to a single
number. (i)~Against the Army--Navy Basic Finner \emph{finned}-vehicle free-flight
$C_{m_q}$ data of Dupuis and Hathaway~\cite{dupuis1997}---a different
configuration and reference from the body-alone Tobak TN~3788 row
(\#16, $39\%$ at $M=1.5$) in Table~\ref{tab:benchmarks}---the augmented model is
sign-correct but exhibits a supersonic MAPE near $69\%$ (under-prediction) over the
supersonic points. (ii)~Against the independent Bhagwandin and Sahu URANS free-flight
comparison~\cite{bhagwandin2013} the same model returns a supersonic MAPE of
$18.96\%$ on the AFF planform (5 points, $M = 1.30$--$2.50$) and $28.02\%$ on the
ANF planform (8 points, $M = 1.29$--$4.50$), sign-consistent with the Basic Finner
under-prediction. (iii)~Against the engineering-model damping CFD of
Sznajder~\cite{sznajder2025}, the transonic Gaussian over-predicts the
free-oscillation damping peak by roughly $+110$ to $+160\%$ at $M = 1.08$--$1.11$;
this overshoot is a comparison versus a different reference dataset and should not
be conflated with the model's internal Gaussian augmentation parameter ($k = 3.5$
at $M = 1$) or with the $\sim 69\%$ Basic Finner free-flight figure. All three are
reported as B-level so downstream users can re-identify the damping closure as
further independent datasets become available.

The hypersonic cone foredrag benchmark, taken against the DTIC AD0487365
cone-cylinder range tests~\cite{grabow1965} with the source Reynolds number
matched row by row, returns a MAPE of $19.7\%$ over eleven points spanning
$M = 6.5$ to $17.2$, with the largest single residual reaching $+57\%$ on the
thinnest, lowest-Reynolds-number cone (Fig.~\ref{fig:cone-drag}). The large residuals
fall on the thinnest ($\theta_c \leq 8^{\circ}$) cones at the lowest-Reynolds-number
rows, where friction and base drag dominate and the AD0487365 boundary-layer
transition state is incompletely specified. This row is therefore an
\emph{exploratory} extension of the supersonic-validated foredrag method and is
reported as B-level rather than as a validated A-level subsystem; until a
thin-cone dataset with a documented transition state is located, the hypersonic
regime as a whole is treated as an exploratory extension rather than a validated
operating range.

Finally, and most importantly for the integrated-flight claim, the finned-vehicle
base-drag augmentation depends on two scale constants that are
\emph{corpus-frozen} rather than externally anchored: a thick-boundary-layer
scale $\mathtt{THICK\_BL\_K} = 2.2$ and a slender-body scale
$\mathtt{SLENDER\_BODY\_K} = 0.0025$. These were fixed against a small set of
amateur flights---the anchor flights Raven, Rabia, Kinsel, and Torrent---so the
flight-corpus apogee result of Sec.~\ref{sec:corpus} is, to that extent,
partly in-sample. We disclose this explicitly and do not present the corpus
aggregate as a fully blind result. The principal defense against an overfitting
interpretation is the decontaminated prospective holdout of
Sec.~\ref{sec:corpus}, in which every flight that either constant touched is
placed in the development split: the genuinely blind holdout is \emph{more}
accurate than the development split, indicating that the two constants generalize
rather than overfit. The reader is referred to Sec.~\ref{sec:corpus} for the
split construction and the corresponding error statistics.

\begin{figure}[tb]
\centering
\includegraphics[width=\columnwidth]{basic_finner_total_drag}
\caption{Finned-vehicle total drag for the Army--Navy Basic Finner reference
projectile versus the Dupuis and Hathaway ballistic-range data~\cite{dupuis1997},
$M = 1.08$--$4.30$. Aggregate MAPE $11.8\%$ over eight points.}
\label{fig:basic-finner}
\end{figure}
